% ===================================================================
%  TAFEL Nr. 9 — Das Gebirge und das Bad. --- Der Wald und die Jagd.
%  Traduction 2026 d'apres texto/io, controlee sur texto/fr et
%  texto/en. Consigne : voir texto/de/10-tafel-01.tex.
%
%  CETTE PAGE EST LA SEULE OU L'IDO ECRIVE UN MOT ALLEMAND SANS SAVOIR
%  QU'IL EST ALLEMAND.
%
%  A l'alinea 16, le narrateur monte « avec son alpenstoko a la
%  main ». Guignon ne le signale pas, l'anglais ecrit « alpenstock »
%  sans y penser, le francais aussi : le mot est passe dans toutes les
%  langues d'Europe sans qu'on s'apercoive qu'il vient de l'allemand
%  des Alpes. Il n'est pas le seul sur cette page. La haute montagne
%  a exporte son vocabulaire entier :
%     Gletscher   (5)   le glacier — du haut-valaisan « glatscher »
%     Lawine            l'avalanche — du romanche « lavina »,
%                       passe par l'allemand de Suisse
%     Alpenstock  (—)   le baton ferre
%     Rucksack    (47)  le sac de course
%     Firn, Gr at, Joch, Kar, Wand : la moitie du vocabulaire
%                       alpiniste de toutes les langues.
%
%  ET C'EST LE CONTRAIRE EXACT DU TABLEAU 2, POUR LA MEME RAISON.
%  Jahn a du IMPOSER « Turnen » a des Allemands qui disaient
%  « Gymnastik », par decret et par campagne. Personne n'a jamais eu a
%  imposer « Gletscher » a personne : le mot est parti tout seul, dans
%  les vingt langues, en un demi-siecle. La difference n'est pas dans
%  la volonte de celui qui nomme, elle est dans la chose :
%
%     UN MOT VOYAGE QUAND LA CHOSE VOYAGE ; IL FAUT LE FORCER
%     SEULEMENT QUAND LA CHOSE EST DEJA LA.
%
%  L'alpinisme est ne vers 1850, chez des gens qui parlaient allemand
%  ou qui apprenaient d'eux ; le gymnase de 1811 se trouvait dans un
%  pays qui avait deja des gymnastes et un mot pour eux. C'est la
%  troisieme fois que la serie mesure cet ecart — apres le calendrier
%  du tableau 3 et l'Achterbahn du meme — et c'est la premiere fois
%  qu'elle le voit du cote de l'exportateur.
%
%  LA SECONDE SCENE DONNE LE QUATRIEME ILOT INTACT. Apres le corps
%  (tableau 2), les metiers (tableau 5) et le betail (tableau 7),
%  voici la foret : Wald, Buche, Birke, Kiefer, Eiche, Efeu, Moos,
%  Specht, Hirsch, Reh, Wildschwein, Hase, Fuchs, Dachs, Förster,
%  Wilderer, Köhler, Holzfäller, Reisig, Rinde, Ast, Zweig, Wurzel.
%  Pas un emprunt sur vingt-trois mots. Les quatre ilots ont ceci de
%  commun qu'on n'y achete rien : le corps, l'atelier, l'etable et la
%  foret sont les quatre endroits ou l'on ne rencontre pas de
%  marchand.
%
%  Une seule exception, et elle confirme : la Gemse de
%  l'alinea 18 et le Murmeltier de l'alinea 15 sont
%  germaniques, mais le Zigeuner de l'alinea 15, qui promene
%  son ours, ne l'est pas — il est grec, par le byzantin. Le seul
%  etranger de la foret est le seul qui n'y habite pas.
% ===================================================================

%%K t09-apar-1 apar
\VUsaut{4.0mm}
\VUtitre{15.0pt}{60}{TAFEL Nr. 9}
\VUsaut{3.0mm}

%%K t09-tit-1 sub
\VUcentre{12.0pt}{0}{\textit{Erste Szene.}}
\VUsaut{3.0mm}

%%K t09-tit-2 sub
\VUpk{11.4pt}{40}{Das Gebirge. --- Der Kurort.}
\VUsaut{3.0mm}

%%K t09-c1-01-1 p
1. --- Ich bin seit einer Woche in Pratz. Ich kann gar nicht sagen, wie sehr ich
gestaunt habe, als ich vor der prächtigen \VUgras{Bergkette} \textsuperscript{(1)} der
Pyrenäen stand, deren \VUgras{Kamm} \textsuperscript{(2)} sich wie eine riesige
Jakobsmuschel vom blauen Himmel abhebt.

%%K t09-c1-02-1 p
2. --- Ich hatte mir nicht vorgestellt, dass die \VUgras{Gipfel}
\textsuperscript{(3)} der Pyrenäen eine solche \VUgras{Höhe} erreichen könnten, und
ich war überwältigt von den prächtigen \VUgras{Gletschern} \textsuperscript{(5)} und
den schneebedeckten \VUgras{Spitzen} \textsuperscript{(4)}, an denen so
schreckliche \VUgras{Lawinen} niedergehen.

%%K t09-tit-3 sub
\VUsaut{3.0mm}
\VUpk{10.2pt}{40}{Gebirgslandschaft.}
\VUsaut{3.0mm}

%%K t09-c1-03-1 p
3. --- Gleich am Tag nach meiner Ankunft bin ich zu dem alten erloschenen
\VUgras{Vulkan} \textsuperscript{(6)} bei Pratz gegangen. Am kahlen, öden Rand des
\VUgras{Kraters} sieht man noch deutlich die Spuren, die die \VUgras{Lava}
hinterlassen hat, die vor einigen Jahrhunderten den \VUgras{Berghang}
\textsuperscript{(7)} hinuntergeflossen ist. An einem der \VUgras{Hänge} habe ich ein
paar Stücke dieser Lava gefunden.

%%K t09-c1-04-1 p
4. --- Pratz liegt an der Grenze, und die \VUgras{Festung} \textsuperscript{(8)}, die
den \VUgras{Pass} \textsuperscript{(27)} zwischen Frankreich und Spanien bewacht,
sieht ganz wie jene alten Burgen am Rhein aus, die von ihrem Felsen herab nach
\VUgras{Beute} auszuspähen scheinen.

%%K t09-c1-05-1 p
5. --- Ein riesiger \VUgras{Adler} \textsuperscript{(9)}, der seinen Horst nicht weit
vom \VUgras{Fort} \textsuperscript{(8)} hat, zieht oft meine Aufmerksamkeit auf sich.
Dieser \VUgras{Raubvogel} mit seinem kräftigen \VUgras{Schnabel} bringt seinen
Jungen oft ein Lamm oder ein Zicklein in den \VUgras{Fängen}. Seine
\VUgras{Flügel} sind so groß, dass er mit seiner schweren Last lange gleiten kann.

%%K t09-tit-4 sub
\VUsaut{3.0mm}
\VUpk{10.2pt}{40}{Der Wanderer.}
\VUsaut{3.0mm}

%%K t09-c1-06-1 p
6. --- Der schönste Spaziergang hier ist der Weg zum \VUgras{Wasserfall}
\textsuperscript{(10)} von Cau. Das \VUgras{stürzende Wasser} fällt wirbelnd herab
und verschwindet dann als hübscher \VUgras{Bach} im \VUgras{Tannenwald}
\textsuperscript{(11)} westlich der Stadt. Wir sind durch diesen Wald bis zur
\VUgras{Wetterwarte} \textsuperscript{(12)} hinaufgestiegen, und von der
\VUgras{Aussichtskanzel} aus hatten wir das ganze \VUgras{Panorama} der Gegend vor
uns. Die \VUgras{Landschaft} ist großartig, und die \VUgras{Natur} scheint alles,
was sie hat, über dieses Land ausgeschüttet zu haben.

%%K t09-c1-07-1 p
7. --- Zum Hinuntergehen haben wir die \VUgras{Bergbahn} \textsuperscript{(13)}
genommen und an einer kleinen \VUgras{Station} bei den \VUgras{Ruinen}
\textsuperscript{(14)} einer alten \VUgras{Ritterburg} gehalten, in der einst die
Herren von Pratz gewohnt haben.

%%K t09-c1-08-1 p
8. --- Arme Burg! Was mag sie sich denken, wenn sie unter sich den hübschen
\VUgras{Badeort} \textsuperscript{(15)} sieht, der jetzt ein modisches \VUgras{Bad}
ist?

%%K t09-c1-08-2 p
Das \VUgras{Klima} ist so angenehm und das \VUgras{Heilwasser} so wohltuend, dass
die \VUgras{Kur} im \VUgras{Sanatorium} \textsuperscript{(17)} ein wahres Vergnügen
ist.

%%K t09-tit-5 sub
\VUsaut{3.0mm}
\VUpk{10.2pt}{40}{Ein hübscher Badeort.}
\VUsaut{3.0mm}

%%K t09-c1-09-1 p
9. --- Man hat wirklich alles getan, um den Aufenthalt angenehm zu machen. Ein
großer \VUgras{Viadukt} \textsuperscript{(16)} wurde für eine Bahnlinie gebaut; der
\VUgras{Park} \textsuperscript{(18)} des \VUgras{Kurhauses}, in dem
\VUgras{Konzerte} gegeben werden, ist reizend angelegt. Mehrere zierliche \VUgras{Villen}
\textsuperscript{(19)} sind rund um das \VUgras{Kasino} \textsuperscript{(21)}
gebaut worden, und eine sehr gut gepflegte \VUgras{Straße} \textsuperscript{(22)}
macht das Fortkommen ausgesprochen leicht.

%%K t09-c1-10-1 p
10. --- Nur ein \VUgras{Damm} \textsuperscript{(23)} trennt den \VUgras{Weg}
\textsuperscript{(20)} vom eigentlichen \VUgras{Tal} \textsuperscript{(25)}, in dessen
Grund ein zierlicher kleiner \VUgras{See} \textsuperscript{(26)} liegt, der einen
klaren \VUgras{Bach} \textsuperscript{(24)} speist.

%%K t09-c1-11-1 p
11. --- Auf der anderen Seite des Tals wird ein \VUgras{Eisenbergwerk}
\textsuperscript{(28)} betrieben. Mehrmals bin ich hingegangen und habe zugesehen,
wie die \VUgras{Bergleute} in den \VUgras{Schacht} einfuhren, und ich bin sogar in
den \VUgras{Stollen} hineingegangen, wo sie den Tag über das \VUgras{Erz} brechen,
das das \VUgras{Metall} enthält.

%%K t09-c1-12-1 p
12. --- Neben dem Bergwerk wurde eine große \VUgras{Hütte} gebaut, um den Reichtum,
den man ihm entnimmt, gleich an Ort und Stelle zu verwerten. Der
\VUgras{Hochofen} \textsuperscript{(29)}, den man mit der hier reichlich
vorhandenen \VUgras{Kohle} heizt, macht es möglich, schnell und billig
\VUgras{Roheisen} zu gewinnen.

%%K t09-c1-13-1 p
13. --- Nicht weit vom Bergwerk wird ein \VUgras{Steinbruch} \textsuperscript{(30)}
abgebaut. Was der alte \VUgras{Steinbrecher} mit seiner \VUgras{Spitzhacke}
losschlägt, ist weder \VUgras{Granit} noch \VUgras{Porphyr} noch auch
\VUgras{Sandstein}, sondern schlichter \VUgras{Bruchstein} zum Häuserbauen.

%%K t09-c1-14-1 p
14. --- Eine der schönsten Zierden dieses Landes ist die \VUgras{Tanne}
\textsuperscript{(31)}. Überall --- am Grund einer \VUgras{Schlucht}
\textsuperscript{(32)}, am Ufer eines \VUgras{Wildbachs} \textsuperscript{(33)}, neben
einem \VUgras{Abgrund} \textsuperscript{(34)} oder einem Steilhang --- setzen ihre
schlanken Nadeln ihren grünen Punkt.

%%K t09-tit-6 sub
\VUsaut{3.0mm}
\VUpk{10.2pt}{40}{Wandern.}
\VUsaut{3.0mm}

%%K t09-c1-15-1 p
15. --- Ich nutze die Ferien für lange Wanderungen. Auf den \VUgras{Pfaden}
\textsuperscript{(35)}, die sich den Berg hinaufwinden, kann man mehrere
\VUgras{Kilometer} und oft mehrere \VUgras{Meilen} zurücklegen, ohne die
Entfernung zu merken.

%%K t09-c1-15-2 p
\VUgras{Meilensteine} \textsuperscript{(36)} und \VUgras{Wegweiser}
\textsuperscript{(37)} bezeichnen die Pfade, auf denen man oft einem jungen
\VUgras{Bergburschen} \textsuperscript{(38)} begegnet, der einen \VUgras{Ranzen}
\textsuperscript{(40)} an der Seite trägt und ein \VUgras{Murmeltier}
\textsuperscript{(39)} dabeihat, oder einem \VUgras{Zigeuner} \textsuperscript{(41)},
dessen \VUgras{Pelzmütze} \textsuperscript{(42)} sonderbar zu seiner zerrissenen
alten \VUgras{Kapuze} \textsuperscript{(43)} passt. Er führt einen dressierten
\VUgras{Bären} \textsuperscript{(44)} an einer \VUgras{Kette}.

%%K t09-tit-7 sub
\VUpk{10.2pt}{40}{Bergsteigen.}
\VUsaut{3.0mm}

%%K t09-c1-16-1 p
16. --- Fast jeden Tag gehe ich mit einem \VUgras{Bergführer} \textsuperscript{(48)}
auf eine \VUgras{Bergtour}, und ich bin sehr stolz, wenn ich mit meinem
\VUgras{Rucksack} \textsuperscript{(47)} auf dem Rücken und meinem
\VUgras{Alpenstock} in der Hand wie ein richtiger \VUgras{Tourist}
\textsuperscript{(46)} die \VUgras{Felsen} \textsuperscript{(45)} angehe.

%%K t09-c1-17-1 p
17. --- Von einem Berggipfel aus sehen die Leute in der Ebene nicht größer aus als
Ameisen; eine \VUgras{Schafherde} \textsuperscript{(50)}, die auf einer
\VUgras{Weide} \textsuperscript{(49)} grast, sieht aus wie ein Kinderspielzeug. Eine
\VUgras{Kiefer} \textsuperscript{(51)} oder auch ein \VUgras{Holzchalet}
\textsuperscript{(52)} ist kaum mehr als ein Fleck im Grün.

%%K t09-c1-18-1 p
18. --- Auf diesen Touren begegnen wir auf dem \VUgras{Steig}
\textsuperscript{(53)} oft einem \VUgras{Gemsjäger} \textsuperscript{(54)}, der das
eben erlegte Tier auf der Schulter trägt.

%%K t09-c1-19-1 p
19. --- In mein Herbarium habe ich einen sehr bemerkenswerten Wedel
\VUgras{Farn} \textsuperscript{(55)} gelegt und einen hübschen rosa Zweig
\VUgras{Heidekraut} \textsuperscript{(57)}, den ich auf meiner letzten Wanderung am
Eingang einer kleinen \VUgras{Höhle} \textsuperscript{(56)} gepflückt habe.

%%K t09-tit-8 sub
\VUsaut{3.0mm}
\VUcentre{12.0pt}{0}{\textit{Zweite Szene.}}
\VUsaut{3.0mm}

%%K t09-tit-9 sub
\VUpk{11.4pt}{40}{Der Wald. --- Die Jagd.}
\VUsaut{3.0mm}

%%K t09-c2-01-1 p
1. --- Gestern habe ich einen herrlichen Tag im Wald verbracht. Das geschäftige
Leben der Tiere und der \VUgras{Insekten} \textsuperscript{(5)}, die dort leben, hat
mich sehr fasziniert.

%%K t09-c2-01-2 p
Die \VUgras{Spinne} \textsuperscript{(1)}, die ihr \VUgras{Netz}
\textsuperscript{(2)} zwischen zwei Zweigen spinnt, um die vorbeifliegende
\VUgras{Fliege} \textsuperscript{(3)} zu fangen, die \VUgras{Schnecke}
\textsuperscript{(4)}, die sich mit ihrem Haus mühsam voranschleppt, der Hirschkäfer
auf der Suche nach Beute, die fleißige Ameise bei der Arbeit neben dem
\VUgras{Ameisenhaufen} \textsuperscript{(6)} --- alle diese kleinen Geschöpfe kamen
und gingen und arbeiteten mit außerordentlicher Kraft.

%%K t09-c2-02-1 p
2. --- Eine \VUgras{Schlange} \textsuperscript{(7)}, die ich zuerst für eine
\VUgras{Viper} hielt, die aber nur eine harmlose \VUgras{Ringelnatter} war, lag
unter einem Baumstamm zusammengerollt und streckte ab und zu ihren schmalen kleinen
Kopf heraus, der immer bereit schien zuzubeißen. Vor mir kroch eine große
\VUgras{Nacktschnecke} \textsuperscript{(8)} schwerfällig dahin, nachdem sie eine der
schmackhaften kleinen \VUgras{Walderdbeeren} \textsuperscript{(11)} verschlungen
hatte, die hier reichlich wachsen.

%%K t09-c2-03-1 p
3. --- Der Boden war mit \VUgras{Waldblumen} bedeckt: \VUgras{Schlüsselblumen}
\textsuperscript{(9)}, \VUgras{Immergrün} \textsuperscript{(10)} und viele andere
Pflanzen, die ich nicht benennen konnte, blühten zwischen den
\VUgras{Grasbüscheln} \textsuperscript{(12)}.

%%K t09-c2-04-1 p
4. --- In dieser Gegend ist der \VUgras{Trüffel} fast so häufig wie der
\VUgras{Pilz} \textsuperscript{(14)}, und ein \VUgras{junges Schwein}
\textsuperscript{(13)}, das einem Förster gehörte, suchte gierig danach, ob es
welche finden könnte.

%%K t09-tit-10 sub
\VUsaut{3.0mm}
\VUpk{10.2pt}{40}{Wilderei.}
\VUsaut{3.0mm}

%%K t09-c2-05-1 p
5. --- Ich habe das \VUgras{Ende} einer Szene gesehen, die ich sehr komisch fand.
Dank der Nase seines \VUgras{Dachshunds} \textsuperscript{(15)} hatte der
\VUgras{Förster} \textsuperscript{(16)} gerade einen \VUgras{Wilderer}
\textsuperscript{(22)} in dem Augenblick überrascht, als der Mann das erlegte
\VUgras{Wild} \textsuperscript{(17)} forttragen wollte. Der Wilderer wollte lieber
seine Beute aufgeben, als verhaftet zu werden, und war geflohen, und
\VUgras{Hase} \textsuperscript{(18)}, \VUgras{Hirschkuh} \textsuperscript{(19)},
\VUgras{Reh} \textsuperscript{(20)} und \VUgras{Wildschwein}
\textsuperscript{(21)} lagen zur Freude des Försters im Gras.

%%K t09-c2-06-1 p
6. --- Die Vielfalt der Bäume, die der Wald enthält, ist erstaunlich. Die
\VUgras{Buchen} \textsuperscript{(23)} und die \VUgras{Birken}
\textsuperscript{(26)}, die \VUgras{Kiefern}, die \VUgras{Harz}
\textsuperscript{(27)} geben, die \VUgras{Eichen} \textsuperscript{(29)} und viele
andere dazu mischen ihre Zweige.

%%K t09-c2-06-2 p
Der \VUgras{Efeu} \textsuperscript{(24)}, der sich an die Stämme der hohen Bäume
klammert, gibt dem \VUgras{Hochwald} \textsuperscript{(25)} ein leuchtendes Grün, das
sich angenehm mit dem blasseren, weicheren Ton des \VUgras{Mooses}
\textsuperscript{(31)} mischt, in dem der \VUgras{Grünspecht} \textsuperscript{(30)}
nach Insekten sucht.

%%K t09-tit-11 sub
\VUsaut{3.0mm}
\VUpk{10.2pt}{40}{Waldlandschaft.}
\VUsaut{3.0mm}

%%K t09-c2-07-1 p
7. --- Die alten Eichen haben mich bezaubert. Ihre dicken, verdrehten
\VUgras{Wurzeln} \textsuperscript{(32)}, halb aus dem Boden heraus, geben ihnen ein
prächtiges Aussehen von Kraft und Saft, und ich frage mich, wie der
\VUgras{Besitzer} \textsuperscript{(35)} es übers Herz bringt, sie fällen und der
\VUgras{Säge} \textsuperscript{(34)} des \VUgras{Holzfällers} \textsuperscript{(33)}
übergeben zu lassen. Die armen \VUgras{Stämme} \textsuperscript{(36)}, ihrer
\VUgras{Rinde} \textsuperscript{(37)} beraubt oder von der \VUgras{Axt}
\textsuperscript{(38)} des \VUgras{Tagelöhners} \textsuperscript{(39)} gespalten,
tun mir leid.

%%K t09-c2-08-1 p
8. --- Ganz in der Nähe hatte ein alter \VUgras{Köhler} \textsuperscript{(41)} mit
seiner \VUgras{Schaufel} einen \VUgras{Meiler} \textsuperscript{(42)} Holzkohle
aufgesetzt, und eine alte Frau trug ein \VUgras{Bündel} \textsuperscript{(40)}
\VUgras{Reisig} fort, das aus den Zweigen der gefällten Bäume gemacht war.

%%K t09-tit-12 sub
\VUsaut{3.0mm}
\VUpk{10.2pt}{40}{Die Jagd.}
\VUsaut{3.0mm}

%%K t09-c2-09-1 p
9. --- Ich wollte mich eine Weile im \VUgras{Forsthaus} \textsuperscript{(46)}
ausruhen, aber als ich an den kleinen \VUgras{Teich} \textsuperscript{(43)} kam, auf
dem ein schöner \VUgras{Schwan} \textsuperscript{(45)} schwamm, merkte ich, dass ein
frisches \VUgras{Hochwasser} \textsuperscript{(44)} den Weg in einen richtigen
\VUgras{Sumpf} verwandelt hatte. Ich musste einen Umweg machen, und als ich an der
\VUgras{Zeder} \textsuperscript{(49)}, den \VUgras{Pappeln} \textsuperscript{(48)} und
der \VUgras{Ulme} \textsuperscript{(47)} am Waldrand vorbeikam, sah ich aus einem
\VUgras{Dickicht} \textsuperscript{(54)} einen prächtigen \VUgras{Hirsch}
\textsuperscript{(50)} brechen, gejagt von einer unerbittlichen \VUgras{Meute}
\textsuperscript{(51)}, die das \VUgras{Horn} \textsuperscript{(53)} der
\VUgras{berittenen Jäger} \textsuperscript{(52)} antrieb. Das arme Tier war völlig
erschöpft. Ein paar Schritte von mir entfernt brach es sterbend zusammen.

%%K t09-c2-10-1 p
10. --- Der Sohn des Försters, den der Lärm der \VUgras{Jagd} herausgelockt hatte,
kam aus dem Haus, und wir gingen zu zweit wieder in den Wald hinein. Er hatte eine
\VUgras{Elster} \textsuperscript{(56)} auf dem Ast einer alten Eiche gesehen. Er
meinte, ihr Nest müsse im \VUgras{Laub} \textsuperscript{(58)} versteckt sein, und
wollte es ausnehmen.

%%K t09-c2-10-2 p
Flink wie ein \VUgras{Eichhörnchen} \textsuperscript{(61)} kletterte er von
\VUgras{Ast} \textsuperscript{(55)} zu Ast, fand aber nichts und scheuchte nur eine
alte \VUgras{Krähe} \textsuperscript{(60)} auf, die auf einem \VUgras{Zweig}
\textsuperscript{(57)} saß.
\VUsaut{4.0mm}
\VUfilet{22mm}
